% =========================================================================
% macros.tex - 全局学术宏包、专用数学符号与图表高亮命令
% =========================================================================

% 基础数学与字体支持
\usepackage{amsmath,amssymb,amsfonts}
\usepackage{bm}                       % 粗体数学符号
\usepackage{graphicx}
% Keep referenced floats within their discussion before the conclusion.
\usepackage{placeins}
\usepackage{booktabs}                 % 标准三线表 (\toprule, \midrule, \bottomrule)
\usepackage{multirow}
\usepackage{makecell}
\usepackage{array}
\usepackage{xcolor}
\usepackage{colortbl}

% 顶刊紧凑版面微调
\let\Bbbk\relax                       % 避免符号冲突
\usepackage{newtxtext,newtxmath}
\usepackage[font=small,labelfont=bf,labelsep=period]{caption}
\usepackage[font=footnotesize]{subcaption}

% 交叉引用与超链接
\usepackage[pagebackref,breaklinks,colorlinks,citecolor=blue,linkcolor=red]{hyperref}
\usepackage[capitalize]{cleveref}
\usepackage{pgfplots}

% 常用学术缩写
\newcommand{\etal}{\textit{et al.}}
\newcommand{\ie}{\textit{i.e.}}
\newcommand{\eg}{\textit{e.g.}}
\newcommand{\wrt}{\textit{w.r.t.}}

% 表格指标高亮
\newcommand{\best}[1]{\textbf{\textcolor{red}{#1}}}
\newcommand{\second}[1]{\underline{\textcolor{blue}{#1}}}

% 图像与空间定义
\newcommand{\R}{\mathbb{R}}
\newcommand{\Iir}{\mathbf{I}_{\text{ir}}}       % 红外源图像
\newcommand{\Ivis}{\mathbf{I}_{\text{vis}}}     % 可见光源图像
\newcommand{\Ifuse}{\mathbf{I}_{\text{f}}}      % 最终融合图像
\newcommand{\Irec}{\hat{\mathbf{I}}}            % Phase I 自编码重构图像

% 多阶段特征张量表达
\newcommand{\Fsh}{\bm{\Phi}^{\text{S}}}         % 浅层共享特征 (Stage 1)
\newcommand{\Fbase}{\bm{\Phi}^{\text{B}}}       % 低频基底特征 (Stage 2A)
\newcommand{\Fdet}{\bm{\Phi}^{\text{D}}}        % 高频细节特征 (Stage 2B)

% 网络组件算子
\newcommand{\Eshared}{\mathcal{E}_{\text{stem}}}
\newcommand{\Ebase}{\mathcal{E}_{\text{B}}}
\newcommand{\Edetail}{\mathcal{E}_{\text{D}}}
\newcommand{\Drec}{\mathcal{D}}
\newcommand{\FfuseB}{\mathcal{F}_{\text{B}}}
\newcommand{\FfuseD}{\mathcal{F}_{\text{D}}}

% 损失函数符号体系
\newcommand{\Ltotal}{\mathcal{L}_{\text{total}}}
\newcommand{\Ldecomp}{\mathcal{L}_{\text{decomp}}^{\text{ortho}}}
\newcommand{\Lintra}{\mathcal{L}_{\text{intra}}}
\newcommand{\Lfusion}{\mathcal{L}_{\text{fusion}}}
\newcommand{\Lin}{\mathcal{L}_{\text{in}}}
\newcommand{\Lgrad}{\mathcal{L}_{\text{grad}}}
\newcommand{\Lstruct}{\mathcal{L}_{\text{struct}}}
\newcommand{\Lsd}{\mathcal{L}_{\text{sd}}}
\newcommand{\CC}{\operatorname{CC}}             % 皮尔逊相关系数

\pgfplotsset{compat=1.18}
